\documentclass{article}
\usepackage{Self_Style}

\title{UChicago REU Notes Week 5}
\author{Zih-Yu Hsieh}
\date{\today}

\begin{document}
\maketitle
\tableofcontents

To be honest, last week's note isn't as successful, so we'll mainly focus on the category theory from now on. Here, we assume that $G$ is a finite group, and the categories (mentioned as objects of $\Cat$ or similar objects) are small.

\section{Statements on $G$-categories}

This part is devoted to understand some basics. I'll fill in the details for preliminaries on $G$-categories later, together with adjunction $(\Ob\dashv \mathcal{E})$ between $\Cat$ and $\Sets$.

Given a category $\mathcal{A}$, the functor category $\Cat(\mathcal{\mathcal{E}(G)},\mathcal{A})$ is automatically a $G$-category, as each element $g\in G$ can be realized as an automorphism functor on the indiscrete category $\mathcal{E}G$, so it admits both a left/right action by precomposing either $g^{-1}$ or $g$ as automorphism of $\mathcal{E}G$.

Similarly, if $\mathcal{A}$ itself already has a left/right $G$-action, then it defines a left/right $G$-action on $\Cat(\mathcal{B},\mathcal{A})$ by postcomposition; similarly if $\mathcal{B}$ is a $G$-category, then precomposition also defines one. Combining the two, it admits a conjugation action on $\Cat(\mathcal{B},\mathcal{A})$ if both $\mathcal{B},\mathcal{A}$ admits a left action.

Given $G$-categories $\mathcal{A},\mathcal{B}$, denote $\Cat_G(\mathcal{B},\mathcal{A}):= \Cat(\mathcal{B},\mathcal{A})$, and $G\Cat(\mathcal{B},\mathcal{A}):=\Cat(\mathcal{B},\mathcal{A})^G$ as the $G$-fixpoint (i.e. all functors commuting with $G$-actions).

\hfil

\begin{prop}
There is a $G$-equivarant isomorphism when giving $\mathcal{A},\mathcal{B},\mathcal{C}$ three $G$-categories:
\[\Phi:\Cat_G(\mathcal{A}\times\mathcal{B},\mathcal{C})\cong \Cat_G(\mathcal{A},\Cat(\mathcal{B},\mathcal{C}))\]
\end{prop}
\begin{proof}
    We recognize this as a non-equivariant statement first, then show if all three are $G$-categories (with $G$ acts diagonally on $\mathcal{A}\times\mathcal{B}$), such isomorphism is $G$-equivariant.

    \hfil

    \textbf{I. Construction of Functor:}

    Given a functor $F:\mathcal{A}\times\mathcal{B}\rightarrow\mathcal{C}$, for any object $X\in\mathcal{A}$ (together with identity $\id_X:X\rightarrow X$), define a functor $F_X:\mathcal{B}\rightarrow \mathcal{C}$ by $F_X(Y):= F(X,Y)$, and given any morphism $g:Y\rightarrow Y'$ in $\mathcal{B}$, define the morphism $F_Xg:F_X(Y)\rightarrow F_X(Y')$ as the morphism $F(\id_X,g):F(X,Y)\rightarrow F(X,Y')$. The routine check that this is functorial, using the product category definition:
    \begin{itemize}
        \item Given $\id_Y:Y\rightarrow Y$, we have $F_X\id_Y = F(\id_X,\id_Y) = \id_{F(X,Y)} = \id_{F_X(Y)}$.
        \item Given $g:Y\rightarrow Y'$, $g':Y'\rightarrow Y''$, we have
        \[F_X(g'\circ g) = F(\id_X,g'\circ g) = F(\id_X,g')\circ F(\id_X,g) = F_Xg'\circ F_Xg\]
    \end{itemize}
    Now, given a morphism $f:X\rightarrow X'$ in $\mathcal{A}$, we shall define a natural transformation $Nf:F_X\rightarrow F_{X'}$. Realistically, for any $Y\in\mathcal{B}$, the only choice is by $Nf_Y:=F(f,\id_Y):F(X,Y)\rightarrow F(X',Y)$. Which, given $g:Y\rightarrow Y'$ in $\mathcal{B}$, the following diagram commutes using product category property:
    \[\xymatrix{
        F(X,Y) \ar[d]_-{F(f,\id_Y)} \ar[r]^-{F(\id_X,g)} & F(X,Y') \ar[d]^-{F(f,\id_{Y'})}\\
        F(X',Y)\ar[r]_-{F(\id_{X'},g)} & F(X',Y')
    }\ \equiv\ \xymatrix{
        F_X(Y) \ar[d]_-{Nf_Y} \ar[r]^-{F_Xg} & F_X(Y') \ar[d]^-{Nf_{Y'}}\\
        F_{X'}(Y)\ar[r]_-{F_{X'}g} & F_{X'}(Y')
    }\]
    Hence, $Nf:F_X\rightarrow F_{X'}$ is a natural transformation.

    Then, we claim that $N:\mathcal{A}\rightarrow \Cat(\mathcal{B},\mathcal{C})$ by $N(X):= F_X$, $Nf:F_X\rightarrow F_{X'}$ as above is a functor:
    \begin{itemize}
        \item $(N\id_X)_Y = F(\id_X,\id_Y) = \id_{F(X,Y)}=\id_{F_X(Y)}$, so $N\id_X=\id_{F_X}$.
        \item Given $f:X\rightarrow X'$, $f':X'\rightarrow X''$, then
        \[N(f'\circ f)_Y = F(f'\circ f,\id_Y) = F(f',\id_Y)\circ F(f,\id_Y) = Nf'_Y\circ Nf_Y\]
        So, $N(f'\circ f)=Nf'\circ Nf$.
    \end{itemize}

    \hfil

    Now, suppose two such functors $F,F':\mathcal{A}\times\mathcal{B}\rightarrow\mathcal{C}$ has a natural transformation $\mu:F\rightarrow F'$, we want to construct a natural transformation $\overline{\mu}:N\rightarrow N'$, where $N,N':\mathcal{A}\rightarrow \Cat(\mathcal{B},\mathcal{C})$ are functors associated to $F,F'$. Which, given any $X\in\mathcal{A}$, we see the collection $(\overline{\mu}_X)_Y:= \mu_{(X,Y)}$ is a natural transformation $F_X\rightarrow F'_X$, due to the following diagram for all $g:Y\rightarrow Y'$ in $\mathcal{B}$:
    \[\xymatrix{
        F(X,Y) \ar[d]_-{\mu_{(X,Y)}} \ar[r]^-{F(\id_X,g)} & F(X,Y') \ar[d]^-{\mu_{(X,Y')}}\\
        F'(X,Y)\ar[r]_-{F'(\id_X,g)} & F'(X,Y')
    }\ \equiv\ \xymatrix{
        F_X(Y) \ar[d]_-{(\overline{\mu}_X)_Y} \ar[r]^-{F_Xg} & F(X,Y') \ar[d]^-{(\overline{\mu}_X)_{Y'}}\\
        F'_X(Y)\ar[r]_-{F'_Xg} & F'_X(Y')
    }\]
    This shows $\overline{\mu}_X:F_X\rightarrow F'_X$ is a natural transformation, or morphism within $\Cat(\mathcal{B},\mathcal{C})$. Then, notice that if $\mu:F\rightarrow F'$, $\mu':F'\rightarrow F''$ are two natural transformations between functors $F,F',F'':\mathcal{A}\times\mathcal{B}\rightarrow\mathcal{C}$, the following holds:
    \[(\overline{\mu'\circ\mu}_X)_Y = (\mu'\circ\mu)_{(X,Y)} = \mu'_{(X,Y)}\circ \mu_{(X,Y)} = (\overline{\mu'}_X)_Y\circ (\overline{\mu}_X)_Y\]
    This shows as natural transformations, $\overline{\mu}_X:N_X\rightarrow N'_X$ and $\overline{\mu'}_X:N'_X\rightarrow N''_X$ satisfies $(\overline{\mu'\circ\mu})_X=\overline{\mu'}_X\circ \overline{\mu}_X$, so $\overline{\mu'\circ\mu} = \overline{\mu'}\circ \overline{\mu}$ as natural transformations $\overline{\mu}:N\rightarrow N'$, $\overline{\mu'}:N'\rightarrow N''$.

    So, define $\Phi(F:\mathcal{A}\times\mathcal{B}\rightarrow\mathcal{C}):= N:\mathcal{A}\rightarrow \Cat(\mathcal{B},\mathcal{C})$, and $\Phi(\mu:F\rightarrow F'):= \overline{\mu}:N\rightarrow N'$ as above, it's a legitimate functor.

    \hfil

    \textbf{II. Isomorphism on Objects:}

    Suppose $\Phi(F)=\Phi(F')$, then it's clear $F,F'$ have the same action on the objects (as all $X\in\mathcal{A}$ satisfies $F_X = F'_X$, so all $Y\in\mathcal{B}$ satisfies $F(X,Y)=F_X(Y)=F'_X(Y)=F'(X,Y)$). Now, given $f:X\rightarrow X'$ in $\mathcal{A}$ and $g:Y\rightarrow Y'$ in $B$, we also see that:
    \[F(f,g)=F(f,\id_{Y'})\circ F(\id_X,g) = Nf_{Y'}\circ F_{X}g = F'f_{Y'}\circ F'_Xg = F'(f,\id_{Y'})\circ F'(\id_X,g) = F'(f,g)\]
    ($Nf=N'f$ is by $\Phi(F)=\Phi(F')$, then $F_X=F'_X$ also the same logic).

    This shows $F=F'$, as its map on morphisms are also the same.

    \hfil

    Conversely, given a functor $G:\mathcal{A}\rightarrow\Cat(\mathcal{B},\mathcal{C})$, we'd like to construct a functor $\tilde{G}:\mathcal{A}\times \mathcal{B}\rightarrow\mathcal{C}$. Object wise, define $G(X,Y):= [G(X)](Y)$ (here, $G(X):\mathcal{B}\rightarrow\mathcal{C}$ is a functor), and given morphisms $f:X\rightarrow X'$ in $\mathcal{A}$, $g:Y\rightarrow Y'$ in $\mathcal{B}$, define morphism $G(f,g):G(X,Y)\rightarrow G(X',Y')$ as the following morphism in the commutative diagram:
    \[\xymatrix{
        G(X,Y):=[G(X)](Y) \ar[d]_-{Gf_Y} \ar[r]^-{G(X)g} & [G(X)](Y') \ar[d]^-{Gf_{Y'}}\\
        [G(X')](Y) \ar[r]_-{G(X')g} & [G(X')](Y')=: G(X',Y')
    }\]
    Here, recall that for each $f:X\rightarrow X'$, $Gf:G(X)\rightarrow G(X')$ is a morphism in $\Cat(\mathcal{B},\mathcal{C})$ (i.e. a natural transformation), so the above diagram holds. 

    Now, to verify functoriality, consider the following:
    \begin{itemize}
        \item Given $\id_X$ in $\mathcal{A}$, $\id_Y$ in $\mathcal{B}$, we get $G(\id_X) = \id_{G(X)}$, so each natural transformation $G(\id_X)_Y = \id_{[G(X)](Y)}$. Similarly, since $G(X)$ is a functor, we get $G(X)\id_Y = \id_{[G(X)](Y)}$. So, the diagram becomes:
        \[\xymatrix{
        \tilde{G}(X,Y):=[G(X)](Y) \ar[d]_-{\id_{[G(X)](Y)}} \ar[r]^-{\id_{[G(X)](Y)}} & [G(X)](Y) \ar[d]^-{\id_{[G(X)](Y)}}\\
        [G(X)](Y) \ar[r]_-{\id_{[G(X)](Y)}} & [G(X)](Y)=: \tilde{G}(X,Y)
    }\]
    Hence, $G(\id_X,\id_Y)$ results in the identity of $G(X,Y)$.

        \item Given $f:X\rightarrow X', f':X'\rightarrow X''$ in $\mathcal{A}$, and $g:Y\rightarrow Y', g':Y'\rightarrow Y''$ in $\mathcal{B}$, the following gigantic diagram holds, due to the functoriality $G(f'\circ f)=Gf'\circ Gf$, and each $G(X)(g'\circ g)=G(X)g'\circ G(X)g$:
        \[\xymatrix{
            \tilde{G}(X,Y):= [G(X)](Y) \ar[r]^-{G(X)g} \ar[d]_-{Gf_Y} & [G(X)](Y') \ar[r]^-{G(X)g'} \ar[d]_-{Gf_{Y'}} & [G(X)](Y'') \ar[d]^-{Gf_{Y''}}\\
            [G(X')](Y) \ar[d]_-{Gf'_Y} \ar[r]^-{G(X')g} & \tilde{G}(X',Y'):= [G(X')](Y') \ar[d]_-{Gf'_{Y'}} \ar[r]^-{G(X')g'} & [G(X')](Y'') \ar[d]^-{Gf'_{Y''}}\\
            [G(X'')](Y) \ar[r]_-{G(X'')g} & [G(X'')](Y') \ar[r]_-{G(X'')g'} & \tilde{G}(X'',Y''):=[G(X'')](Y'')
        }\]
        Here, adjoining diagonal squares provide the composition $\tilde{G}(f',g')\circ \tilde{G}(f,g)$, while the edges of the largest square provides $\tilde{G}(f'\circ f, g'\circ g)$, the commutation shows the functoriality.
    \end{itemize}
    So, by routine check on the definition, we have $\Phi(\tilde{G})=G$, showing surjectivity on the objects.

    \hfil

    \textbf{III. Naturality:}

    Let's skip this...
    

    \hfil

    \textbf{IV. Equivariance Version:}

    Assuming the first three statements, the remaining statement is to show: $\Phi$ as a functor commutes with $G$-action. Let $g\in G$, which given any $F\in\Cat(\mathcal{A}\times\mathcal{B},\mathcal{C})$, we have $g\cdot F:=g\circ F\circ g^{-1}$ being the conjugation action of $G$ on this functor category (where the left $g$ is the functor on $\mathcal{C}$, and the right $g^{-1}$ is the diagonal $g^{-1}$ functor on $\mathcal{A}\times\mathcal{B}$). Let $N := \Phi(F)$, and $N':=\Phi(g\cdot F)$. Then, notice $g\cdot F$ satisfies the following, for all $f:X\rightarrow X'$ in $\mathcal{A}$, and $g:Y\rightarrow Y'$ in $\mathcal{B}$:
    \[N'_X (Y) = g\circ F\circ g^{-1}(X,Y),\quad (g\cdot N)_X (Y) = (g\circ N\circ g^{-1})_X(Y) = (g\circ N_{g^{-1}(X)})(Y)\]
    \[ = g\circ N_{g^{-1}(X)}(g^{-1}(Y)) = g\circ F(g^{-1}(X),g^{-1}(Y))) = g\circ F\circ g^{-1}(X,Y)\]

    (Yeah, we also need to verify that they're identical on the morphisms of $\mathcal{B}$, but let's skip it here...)

    Here, we do need to be careful on if $g$ is a functor on the categories, or the functor on the functor categories (as the second $g\circ N$ is really $g$ acting on the pointwise functor $N_{g^{-1(X)}}$ using conjugation).
    
    This shows $N', g\cdot N$ are identical on the objects. Tp
\end{proof}

God the proof is actually long as hell\dots

\hfil

\begin{defn}
    Given a space $X$, the chaotic topological category $\mathcal{E}X$ has object space $X$, morphism space $X\times X$, such that:
    \begin{itemize}
        \item Identity map $i:X\rightarrow X\times X$ by $x\mapsto (x,x)$ is continuous (also the same as diagonal morphism).
        \item The source map $s:X\times X\rightarrow X$ by $s(x,y):= x$ is continuous (also the same as first projection).
        \item The target map $t:X\times X\rightarrow X$ by $t(x,y):= y$ is continuous (also the same as second projection).
        \item Composition $\circ:(X\times X)\times_{s,t}(X\times X)\rightarrow X\times X$ by $\circ((x,y),(y,z)):= (x,z)$ is continuous.
    \end{itemize}
\end{defn}
Since there is only one morphism between any two objects, a morphism is equivalent to specify its source and target.

In particular, for any $x\in X$, the maps $\eta_x,\mu_x:X\rightarrow X\times X$ by $\eta_x (y):= (x,y)$, and $\mu_x(y) := (y,x)$ are continuous (as they're fixing one coordinate, while doing inclusion on the right). These also form natural isomorphisms $\eta:x\rightarrow X$, and $\mu:X\rightarrow x$ as follow:
\[\xymatrix{
    x \ar[d]_-{\eta_x (y) = (x,y)} \ar[r]^-{(x,x)} & x \ar[d]^-{\eta_x (z)=(x,z)}\\
    y \ar[r]_-{(y,z)} & z
},\quad\quad\xymatrix{
    y \ar[d]_-{\mu_x(y)=(y,x)} \ar[r]^-{(y,z)} & z\ar[d]^-{\mu_x(z)=(z,x)}\\
    x \ar[r]_-{(x,x)} & x
}\]
This shows as a category, $\mathcal{E}X$ is equivalent to the terminal category $*$ (which only has one unique topological category structure).

Also, if $X=G$ is a topological group (including discrete ones!) then $\mathcal{E}G$ is isomorphic to the translation category of $G$ (as all morphisms can be interpreted as a translation from one group element to the other).

\hfil

Let $\Pi$ be a topological group (intepreted as the category of single object $*$, with morphism space $\Pi$), then there is a group homomotphism $\Pi^\Op\rightarrow \Aut_{\mathcal{T}\Cat}(\mathcal{E}\Pi)$, by right multiplication. So, all $g\in \Pi$ satisfies $g:\mathcal{E}\Pi\rightarrow\mathcal{E}\Pi$ by $\sigma\mapsto \sigma g$ (which is a homeomorphism on $\Pi$ by definition of topological group), and morphism $(\sigma,\tau)\mapsto (\sigma g,\tau g)$ (again, this is the continuous diagonal action on $\Pi\times\Pi = \Mor(\mathcal{E}\Pi)$).

\begin{prop}
    The orbit category $\mathcal{E}\Pi/\Pi$ is isomorphic to $\Pi$ as topological category.
\end{prop}
\begin{proof}
    It's clear on the object space, as $\Ob(\mathcal{E}\Pi)=\Pi$ as set, so $\Ob(\mathcal{E}\Pi)\/\Pi = *$.

    The map on the morphism space is $(\sigma,\tau)\sim (\sigma',\tau')$ iff there exists $g\in \Pi$, such that $(\sigma,\tau)=(\sigma'g,\tau'g)$. Which, this is equivalent to say $\sigma\tau^{-1} = \sigma'g g^{-1}\tau'^{-1}= \sigma'\tau'^{-1}$. So, we see all $(\sigma,\tau)\sim (e,\tau\sigma^{-1})$, and all $g,g'\in\Pi$ satisfies $(e,g)\sim(e,g')$ iff $g^{-1}=g'^{-1}$, or $g=g'$. So, we see each $(e,g)$ establishes a distinct equivalent class, hence providing a bijection between $(\Pi\times \Pi)/\Pi\cong\Pi$, by $(e,g)\mapsto g$.

    Finally, to show this is an isomorphism as topological category, we see such identification of morphism is topologically isomorphic (because $\Pi\rightarrow \Pi\times\Pi \rightarrow \Pi$ by $g\mapsto (e,g)\mapsto g$ is a homeomorphism). Then, given $(e,g)$ and $(e,g')$ in the quotient, since the second one is equivalent to $(g,g'g)$, then we see the composition is $\circ((e,g),(g,g'g)) = (e,g'g)$, which has quotient $g'g$, which is equivalent to the composition of the image of $(e,g')$ and $(e,g)$, showing it's preserving the composition map. Hence, it's an isomorphism of categories.
\end{proof}

\hfil

\begin{thm}
Given a $G$-space $X$, a topological group $\Pi$ (regarded as a trivial $G$-space, by $G\rightarrow e\rightarrow \Aut(\Pi)$), the quotient functor $p:\mathcal{E}\Pi\rightarrow\Pi$ induces an isomorphism if topological $G$-categories 
\[\Cat(\mathcal{E}X,\mathcal{E}\Pi)/\Pi\xrightarrow{\sim}\Cat_G(\mathcal{E}X,\Pi)\]
\end{thm}
\begin{proof}
    In terms of objects, we have $\Cat(\mathcal{E}X,\mathcal{E}\Pi)\cong \Sets(X,\Pi)$
\end{proof}

\end{document}